\chapter{基于视频扩散先验的物理参数反演与连续体仿真}

\section{引言}
经过前两章的三维重建与语义特征蒸馏，系统已经成功构建了一个包含高纯净度几何实体与基础物理本构空间的复杂场景。然而，仅凭借泛化视觉大模型给出的粗略材质推断，难以在极其复杂的物理碰撞、挤压或是大变形边界下展现出严丝合缝的力学行为。视觉表象如何向高保真的动力学特性靠拢，是实现"以假乱真"交互的关键。
为此，本章以物质点法（Material Point Method, MPM）作为底层连续体力学仿真底座，首次引入了图像级视频扩散大模型（Stable Video Diffusion, SVD）作为高级物理规律的监督先验。本章将详细论述整个基于视频扩散模型引导的"梦幻物理（DreamPhysics）"寻优求解框架。其核心思想在于：利用扩散网络对自然界物体运动规律的内含认知，引导底层 MPM 引擎在给定的基础本构域（杨氏模量、泊松比）内展开微调梯度搜索。为了解决大尺度场景下的计算灾难问题，本章还专门构建了背景冻结机制与双 GPU 联邦解算架构。

\section{基于物质点法（MPM）的三维连续体力学基础}
鉴于真实世界的物理交互不仅包含刚体位移，往往更伴随着拓扑撕裂、弹性形变、流体飞溅等现象。为了应对这种极其多变的物理状态，本系统集成了物质点法（MPM）作为高保真的仿真底座。

\subsection{质点填充与混合欧拉-拉格朗日求解管线}
不同于三维高斯溅射仅提供一个空间视觉包围盒壳体，物理仿真需要环境拥有真实的"质量体积"。因此，在导入从 Chapter 4 蒸馏输出的独立前景语义后，第一步即是进行内部边界的质点填充（Particle Filling）。每个被网格化（Grid）包裹的有效内部节点都被赋予初始的微观物理参数。
MPM 方法通过在拉格朗日质点（承载质量、动量与形变梯度）和固定的欧拉背景网格（用于高效求解内部非线性力与碰撞边界）之间反复转移张量来求解动量守恒方程。这种混合式架构既有效追踪被摄物体的连续表面演化，又能避免产生纯网格法中常见的大变形网格纠缠畸变，非常契合从粗糙点云初始化的高斯球。

\subsection{混合多材质本构交互与背景冻结机制}
在本系统中，所有交互的拉格朗日质点最初都承接了来源于多模态经验映射得出的 10 大类别初始参量体系。但在室内或室外宏大场景中，若放任整个场景空间全部参与矩阵求解，其庞大而僵硬的背景质量势必如"锚点"一般将整个弹性系统钉死，甚至带来数十倍的无效计算耗时。
以此，我们在仿真引擎前端强制引入了隔离剔除机制（Mark-Foreground）。该系统现已升级为具备自适应预设（compact / wide / multicomponent / full\_object）与多重安全保护的多策略分割框架，其核心采用五阶段精细化管线：

\textbf{阶段一（粗筛选）}：从全体高斯中采样子集，计算各高斯 CLIP 特征与前景文本的余弦相似度，经阈值初筛后进行\textbf{得分引导的 DBSCAN 聚类（Score-Guided DBSCAN）}——不同于传统取最大连通簇的做法，本系统按文本置信度（70\%）、簇尺寸（20\%）与密度（10\%）的加权综合评分选取最优聚类，构建前景粗包围盒并剔除盒外背景。

\textbf{阶段二（精细相似度+实例聚类）}：在粗包围盒内对全量高斯计算多层级相似度，经阈值过滤得到种子掩码，可选地对种子进行 DBSCAN 实例选择与种子包围盒扩展。

\textbf{阶段三（竞争性文本评分与多源证据融合）}：本阶段是系统的核心创新。当同时提供前景与背景文本查询时，系统对所有查询嵌入执行带温度缩放的 softmax 归一化，取前景查询概率之和作为竞争性分数。随后融合四类独立证据源：（1）对象级 CLIP 相似度（权重 0.55）；（2）零件级 CLIP 相似度（权重 0.20）；（3）SAM 二维掩码反投影投票——将 3D 种子投至各视角，匹配 SAM 实例掩码后回投得票（权重 0.20）；（4）DINOv2 特征空间 KNN 平滑余弦相似度（权重 0.05）。这一多源证据融合机制使系统能够在仅靠单一 CLIP 得分无法区分前景与相似背景的困难场景中保持稳健。

\textbf{阶段四（KNN 平滑与连通分量过滤）}：对融合得分进行阈值化后，通过 FAISS-GPU 加速的 KNN 标签传播进行前景膨胀与背景清理。随后利用 GPU 加速的连通分量分析（标签传播最大 64 次迭代的 scatter\_reduce 最小标号传播）提取最大空间连通簇，有效剔除孤立的噪声碎片。

\textbf{阶段五（安全保护与自动回退）}：系统内建多重安全护栏——当前景粒子数超过上限（max\_foreground\_particles=500,000）或前景比例超过预设阈值（compact 模式默认 50\%）时自动收紧阈值；若连通分量过滤导致前景被过度剔除（低于预设最低保留比），则自动回退至过滤前掩码。

只有通过上述五阶段管线确认属于"交互主体"对象的目标团块才会被解冻并激活参与 MPM 形变更新；诸如远山、大地、外壳墙壁皆被声明为静力学无限刚性约束面。此外，针对 CLIP 文本匹配在部分场景（如帽子）中前景召回不佳的问题，系统额外实现了基于 SAM 二维掩码的备选前景标记工具，通过多视角掩码投票将二维分割结果反投影至三维高斯空间，作为文本驱动方法的补充链路。这一机制去除了杂乱背景引入的干扰变量，确保 MPM 求解器只针对真正发生位移的结构对象进行本构演变。

\subsection{语义标签到 MPM 求解器类型的物理参数映射}
前景背景分离完成后，系统需将第四章蒸馏得到的逐高斯语义标签转换为 MPM 仿真引擎可识别的物理参数与求解器类型。该映射由 \texttt{prepare\_dreamphysics\_prerequisites.py} 模块自动完成，其流程包含四项核心安全设计：

\textbf{求解器类型映射（Solver Type Mapping）}：系统维护了一套标签到 MPM 求解器类型的硬编码映射表——例如，金属类（label 1）映射为 von Mises 塑性求解器（type 2），橡胶类映射为 Maxwell 粘弹性求解器（type 9），沙土类映射为 Drucker-Prager 颗粒求解器（type 4）——使得同一场景内的不同材质区域在仿真中使用根本不同的本构模型。

\textbf{少数类中值钳位（Minority Outlier Clamping）}：若某材质类别占激活前景比例低于 1\%，其粒子的杨氏模量与密度被自动钳位至全局活跃中值。这一步防止极个别被误分类为刚性材质（如单个"金属"噪点）的粒子驱动过高的波速，导致仿真时间步长塌陷至不稳定的极小值。

\textbf{基于 CFL 条件的硬 E 上限（CFL-Based Hard E Cap）}：系统从基础 CFL 条件推导各粒子的杨氏模量安全上限：$E_{\max} = v_{\text{wave}}^2 \cdot \rho / \text{wave\_factor}$，其中波速 $v_{\text{wave}} = \mathrm{d}x \cdot 0.1 / \mathrm{d}t$（CFL 安全系数 = 0.1），wave\_factor 通过 $\lambda + 2\mu$ 在泊松比极高时对近不可压缩材料的波速放大进行补偿。任何粒子的 $E$ 超过此上限将被硬截断。

\textbf{自适应网格尺寸（Adaptive Grid Sizing）}：系统根据激活前景点数量自动选取 MPM 背景网格分辨率——20 万粒子以下采用 $\mathrm{d}x = 0.04$，以上放宽至 $0.05$——在仿真精度与显存安全之间取得动态平衡。

\section{基于 Stable Video Diffusion 的时序物理梯度寻优}
即使存在多模态参数的预赋权，该初值依旧属于宏观的"经验范畴"。在不同视角、不同光照条件下，物理实体由于细微瑕疵往往展现出非线性差异。因此，我们利用对自然界视频大数据的生成式规律来隐式地反推与修正物理参数，即 DreamPhysics 框架的时序对齐。

\subsection{选择 SVD 而非传统文本扩散的原因}
在选择视频扩散先验时，我们放弃了传统的基于文本到视频（Text-to-Video, 如 ModelScope）的开源通道。由于在此类任务中，文字不仅无法精确钳制首帧初始态的三维一致性，甚至容易偏离物理常识。相反，Stable Video Diffusion（SVD）属于极其严格的 Image-to-Video（I2V）拓扑结构。它能稳固地锁定场景开始仿真的第零帧图像作为扩散原点，极大程度上维持了高斯空间在静止帧渲染时的极高保真度；且 SVD 具备良好的帧间光流内聚性，十分适用于表征极其复杂的流变或弹性物理下落。

\subsection{反向物理参数寻优与分数蒸馏微分}
DreamPhysics 的核心算法回路是：首发状态的 3DGS-MPM 系统在当前杨氏模量及泊松比下输出一段仿真视频，随后将各帧图序列打包，加注部分环境白噪声后输送至 SVD 网络。通过分数蒸馏采样（Score Distillation Sampling, SDS），系统计算出 SVD 预期未来运动视频与当前 MPM 实际生成视频在微小潜空间内的降噪误差流形。
这一由时序扩散模型产生的分布偏差项不直接作用于图像像素，而是经由基于 Warp 库搭建的全可微（Differentiable）物理管线进行链式法则的复合求导，一路逆向渗透回力学引擎的底层设定，即弹性模量 $E$ 和密度 $\rho$。受梯度驱动，MPM 质点不断调整其内应力对抗，直到其生成的渲染下落轨迹在 SVD 网络看来趋于"天衣无缝的自然"。

为实现上述闭环在高分辨率复杂场景下的稳定运行，系统在仿真运行时引入了五项关键安全保障机制：

\textbf{运行时 CFL 自适应时间步长}：每次执行 MPM 步进前，系统根据当前粒子群的逐点杨氏模量、泊松比和密度重新计算实际波速（以 $\lambda + 2\mu$ 而非简化的 $E/\rho$ 表征近不可压缩材料的波速放大），自适应调整 sub\_step\_dt 与每帧子步数，确保每个渲染帧严格对应一个 frame\_dt 的物理时间推进。

\textbf{逐粒子 CFL 钳位}：SDS 梯度更新后，并非对所有粒子施加统一的 E 上限，而是依据各粒子自身的密度和泊松比逐一计算安全上限。低密度或高泊松比的粒子获得更低的 E 钳位值，有效防止优化器在极个别脆性粒子上注入过大刚度导致仿真发散。

\textbf{体素材质正则化}：在 MPM 网格转移前，系统对粒子材质标签进行体素级多数投票平滑——每个体素内的所有粒子统一采用该体素的大众材质类别，并同步修正其 E、$\nu$、密度与求解器类型。该预处理消除了因孤立噪点粒子引起的局部刚度不连续。

\textbf{按材质类型梯度归一化}：反向传播后，E 的梯度在各材质类型内部独立执行 min-max 归一化至 $[-0.5, 0.5]$ 区间，避免某一材质因参数尺度差异而在梯度上碾压其他材质，确保多材质耦合优化的均衡性。

\textbf{多种振动扰动注入模式}：系统支持四种独立的相机/场景振动模式以激发不同运动模态——solver\_substep 模式在 MPM 每个子步内叠加速度覆盖（产生弯曲变形）、post\_g2p 模式在 G2P 后施加位置修正、force\_substep 模式依中值质量施加力驱动振动、frame\_level 模式在每渲染帧施加刚体速度偏移——用户可根据目标场景的运动特征灵活组合。

\subsection{双 GPU 联邦解算与渲染架构}
在实际工程落地时，MPM 需要进行密集的全可微网格形变张量链式追踪，而 SVD 的批量反向传播极其消耗显存。为此，PhysLucid 提出了针对双卡环境分离处理结构（Dual GPU Layout）：其中 "Render GPU" 承载 MPM 与高斯溅射的混合矩阵计算，专门解算形变梯度的导数及二维像素投影；"Guidance GPU" 则专属常驻 SVD 生成模型与扩散梯度求导网络。两块显卡通过内存切片互锁进行 P2P 梯度矩阵移交。这种联邦工作图谱彻底解放了物理模拟在单卡显存受限时经常遭遇的爆显存（OOM）断点危机，从而保证复杂三维动力现场在高分辨率大算力量级下的反向寻优得以可行。

\section{本章小结}
本章全面阐述了使虚拟视觉环境具备动态逼真物理响应的闭环机制。在物理初始化阶段，通过求解器类型映射、少数类中值钳位、CFL 硬 E 上限与自适应网格尺寸四重安全保障，系统将语义标签稳健地转化为多材质 MPM 求解器配置。在仿真进程中，五阶段多源证据融合的前景分割管线（含得分引导 DBSCAN、竞争性文本评分与四证据源加权融合）配合 SAM 二维掩码反投影补充链路，实现了对复杂场景前景的精准隔离。在 SVD 梯度优化阶段，运行时 CFL 自适应时间步长、逐粒子 CFL 钳位、体素材质正则化与按材质类型梯度归一化等机制共同保障了多材质耦合物理反演在复杂场景下的数值稳定性。这种以扩散模型强大先验为教师、以全可微物理管线为学习终端的梯度反向蒸馏手段，成功消弭了由多模态预估参数造成的呆板视觉差异，并首次突破性地依靠双 GPU 算力层叠实现了复杂材质场景中的微观本构建模。
